\chapter{Introduction}
\label{chap:intro}

In this section we describe the setting of our thesis; we give a detailed
description of the projects we will be working on, the tools we are going to
use, as well as previous research efforts on the area of compiler verification.
We then motivate our project by examining an example that explains the scope of
our work in the PureCake compiler. Finally, we give a brief outline of our work
and the structure of the rest of the thesis.

\section{Formal Verification in Software Systems}

Formal verification is defined as the use of mathematical techniques to ensure
that a design conforms to some precisely expressed notion of functional
correctness \cite{bjesse2005formal}. Specifically, in the context of software
development, a piece of software is said to be formally verified if it provably
transforms its inputs to outputs based on a given specification. A specification
is usually a mathematical model of the program. For example, one could formally
verify a regular expression matching program by describing a mapping of the
program's datatypes to finite state automata, and a mapping of the program's
procedures to transformations over these automata. The specification then can
use the mathematical theory of finite automata to prove theorems about the
program, such as that the matcher only accepts inputs included in the language
of the regular expression.

However, for sufficiently complex programs, doing this transformation on paper
would be impractical and very error-prone. Therefore, the preferred way is to
use software tools called theorem provers or proof assistants. These tools allow
the programmer to state theorems about the behaviour of their code directly in
the programming language and then interactively prove them. In our work for this
thesis, we have used the HOL theorem prover \cite{HOL4}, a proof assistant
implemented on top of the Standard ML language which supports specification and
proof development in higher-order logic.

This thesis is concerned with formal verification of programming language
compilers. Compilers are some of the most complicated pieces of software, and
are notoriously difficult to implement. However, their structure makes them a
particularly good fit to develop in a theorem prover, which presents its own
difficulties, but can guarantee that the compiler is sound and bug-free.
Formally verified compilers already exist and have been successfully used in
real-world systems, where code correctness is critical. A prominent example of a
practical verified compiler is CompCert \cite{CompCert}, a fully-featured,
formally validated compiler for a large subset of the C99 language.

\section{The PureCake Project}

While verified compilers such as CompCert exist, they compile small or low level
languages. This is the case because these languages have been historically used
in critical systems due to their time, memory, and energy efficiency.
Nevertheless, with the constant progress in hardware and software research,
higher-level languages are now also viable to be used in even the most
constrained environments. This is preferred when possible, as high level
languages provide better programming facilities and more guarantees about the
behaviour of the code.

In contrast to low level and usually imperative languages such as C, high level
languages strive to be more declarative and aid the programmer to only have to
think about the task at hand. This is achieved by providing powerful features
both during the compilation of the programs, such as type inference and pattern
matching, and during the runtime such as garbage collection. One such class of
languages are functional languages. Functional programming's main feature is the
treatment of functions as first-class entities, meaning that any function can be
passed as an argument, returned from another function or stored in a data
structure. Whereas in imperative programming procedures are sequences of
instructions, in functional programming functions are expressions that transform
inputs to outputs in a modular manner.

Writing programs in a functional style leads to fewer bugs and easier to test
programs, so having a verified compiler for a functional language would provide
strong guarantees about a program's behaviour. The CakeML project \cite{CakeML}
implements the first fully-featured formally verified compiler for a high level,
functional programming language. It compiles an ML-like language complete with
type inference, pattern matching, and garbage collected runtime, targeting
various hardware architectures. CakeML also implements a series of optimisations
that make the generated executables competitive in runtime efficiency with the
mainstream unverified alternatives.

Even though CakeML is a step towards a more modular and high level verified
language, it follows the tradition of the Standard ML language; code usually
conforms to the functional programming style, but side effects such as mutating
variables and printing to the terminal are still allowed. Pushing the functional
paradigm even further are languages such as Haskell, commonly described as pure,
lazy functional languages. In this context, pure means that calling a function
with the same inputs will always yield the same results, thus implying absence
of side effects, and lazy means that every expression is only evaluated when
needed. This means that computation in these languages directly maps to
evaluation of mathematical functions, which makes reasoning about programs very
intuitive.

Building on top of CakeML, the PureCake project \cite{PureCake} has much the
same goals as CakeML, but for the compilation of lazy pure functional languages.
The PureCake compiler handles the front-end of the compilation process, such as
parsing and type checking. It then transforms the source program to a CakeML
program and hands off the rest of the compilation process to the CakeML
compiler, therefore creating a verified end-to-end compilation process.

\section{PureCake's Compilation Process}

For PureCake to use the CakeML compiler as a backend to handle the verified
generation of machine code, it first has to translate its language to CakeML's
language. The semantics of these languages differ mainly in two ways: purity and
evaluation strategy. Our work is focused on the transformation of the evaluation
strategy, so we won't describe how purity is removed.

CakeML uses the call-by-value evaluation strategy, therefore CakeML is called a
strict language. In a strict language, when a function is called its arguments
are evaluated before the body of the function executes. On the other hand,
PureCake uses the call-by-name evaluation strategy, making PureCake a lazy
language. In PureCake, when a function is called its arguments are not
immediately evaluated. Instead, the function body is evaluated and, if the value
of some argument is demanded for the evaluation of the function body, only then
is the value of the argument computed.

To transition from one evaluation strategy to the other, the PureCake compiler
transforms the lazy program to a strict one by introducing thunks. Thunks are
containers that wrap a single unevaluated expression, and evaluate their
expression when forced to do so. For example, consider the following Pure
function definition \hsinl{suc n = n + 1}. When a call like \hsinl{suc (f x)} is
evaluated, the \hsinl{f x} call is only evaluated when the \hsinl{+} operator
demands its value in the body of \hsinl{suc}. In order to transform the
program's evaluation strategy, the compiler extends the language with two
special expressions: \texttt{delay} and \texttt{force}. \texttt{delay} wraps an
expression in a thunk, and \texttt{force} queries a thunk for its value.
Therefore, the \hsinl{suc (f x)} lazy call could be simulated in a strict
language by transforming the call to \hsinl{suc (delay (f x))} and the function
to \hsinl{suc n = (force n) + 1}.

PureCake achieves the evaluation strategy transformation by passing the source
program through a series of intermediate languages, at each step changing the
semantics of the language to close the gap between PureCake and CakeML's
semantics. The three intermediate languages are \ThunkLang{}, \EnvLang{}, and
\StateLang{}.

\begin{figure}[h]
  \begin{tikzpicture}
    \tikzstyle{box} = [
      rectangle, rounded corners, minimum width=2cm, minimum height=1cm,
      text centered, draw=black
    ]
    \node (box1) [box] {
      \shortstack{PureCake \\ pure \\ call-by-name \\ subst. semantics}
    };
    \node (box2) [box, right=0.85cm of box1] {
      \shortstack {
        \texttt{ThunkLang} \\ pure \\ call-by-value \\ subst. semantics
      }
    };
    \node (box3) [box, right=0.85cm of box2] {
      \shortstack{\texttt{EnvLang} \\ pure \\ call-by-value \\ env. semantics}
    };
    \node (box4) [box, right=0.85cm of box3] {
      \shortstack {
        \texttt{StateLang} \\ impure \\ call-by-value \\ env. semantics
      }
    };
    \draw[-stealth] (box1.east) -- (box2.west) node[]{};
    \draw[-stealth] (box2.east) -- (box3.west) node[]{};
    \draw[-stealth] (box3.east) -- (box4.west) node[]{};
  \end{tikzpicture}
  \caption{PureCake Intermediate Representations}
\end{figure}

In \ThunkLang{}, the \texttt{delay} and \texttt{force} operators are introduced
and placed around the appropriate expressions as described above, and the
evaluation strategy is converted to call-by-value. While this is semantically
enough, it would be inefficient in practice. In \ThunkLang{}, when a thunk is
forced multiple times, its contained expression will be evaluated for each
\texttt{force} separately. Since PureCake is a pure language, we know that the
value of the thunk will always be the same regardless of when we evaluate it.
So, ideally, we should only compute it once and then memoize its value for the
rest of the forces.

This problem is solved in \StateLang{}. \StateLang{}'s semantics employ a
mutable memory called a store. Each time a \texttt{delay} is executed, a thunk
is produced and appended to the store along with a flag stating that the thunk
has not yet been evaluated. Then, the first \texttt{force} demanding the thunk's
value will get a reference to the thunk in the store and evaluate its
expression. It will then set the flag to state that the thunk has been
evaluated, and overwrite the unevaluated expression with the result it computed.
Any subsequent \texttt{force} will then directly use the value instead of
re-computing the expression.

\section{Motivation}

Even though lazy languages provide the facilities to write very high-level code,
laziness comes at a runtime cost. In a strict conventional language, a function
call usually just involves a few machine instructions to move the arguments to
the right registers and push a stack frame. In a lazy language however, we need
to allocate space in the main memory to store a thunk for each argument, which
is orders of magnitudes slower than the strict function calls. Since functions
are invoked constantly in a functional language, it is critical to optimise
their usage as much as possible.

As a case study, consider the following PureCake function that sums the numbers
from 1 to \texttt{n} using an accumulator argument:

\begin{minted}{haskell}
count n acc =
  if n == 0 then acc else count (n - 1) (n + acc)
\end{minted}

We observe that the values of both arguments are always demanded; the value of
\texttt{n} is needed for the conditional, and the value of \texttt{acc} is
demanded either directly in the \texttt{then} case, or through the application
to \hsinl{+} in the \texttt{else} case (this can be inferred by the compiler
using strictness analysis \cite{sergey2014theory}). Therefore, even though this
function is used in a lazy language, its behaviour will always be strict. We
also note that a call-by-value language will optimise this function even further
by noticing it's tail recursive \cite{tailrec}, turning it into a single loop
and thus avoiding even the stack allocations.

Since \hsinl{count} is effectively strict in PureCake too, the compiler should
ideally achieve the same degree of optimisation. Currently, the compiler
initially translates the source language to \ThunkLang{}, by adding delays to
arguments in each function application and forces to every use of a parameter:

\begin{minted}{haskell}
count n acc =
  if (force n) == 0
    then (force acc)
    else count (delay ((force n) - 1))
               (delay ((force n) + (force acc)))
\end{minted}

It then applies some \ThunkLang{}-to-\ThunkLang{} optimisations, namely demand
analysis and a worker-wrapper transformation \cite{Worker/Wrapper}, and produces
the following code:

\begin{minted}{haskell}
count' n acc =
  let n' = force n in
  let acc' = force acc in
  if n' == 0 then acc' else
    count' (delay (n' - 1)) (delay (n' + acc'))
count = delay count'
\end{minted}

The compiler detects that \texttt{n} and \texttt{acc} will always be demanded,
so it forces them in the beginning of the function, thus avoiding work
duplication. Still, each recursive call has to allocate unnecessary thunks for
its arguments.

Our goal in this project is to optimise PureCake's handling of thunks as much as
possible, since it's a critical point for bringing the efficiency of PureCake's
generated code up to par with that of non-verified compilers for lazy languages.

\section{Project Outline}

Our project consists of three parts:
\begin{enumerate}
  \item Creating and formalising a logical relation for \ThunkLang{}
  \item Implementing more \ThunkLang{}-to-\ThunkLang{} optimisations around
    thunks
  \item Adding thunks as first-class constructs in CakeML and optimising their
    representation in the runtime system.
\end{enumerate}

The first part is mostly a preparation for the second part. As it will be
explained in more detail in \autoref{chap:logrel}, developing a logical relation
for \ThunkLang{} will aid in developing and proving optimisations inside
\ThunkLang{}. Furthermore, it can also serve as a general framework for proofs
in \ThunkLang{} that can help reduce boilerplate code and increase proof script
quality and maintainability.

The second part in \autoref{chap:opt} is concerned with developing optimisations
for the handling of thunks in \ThunkLang{}. Going back to the example given in
the previous section, we see that PureCake doesn't currently optimise the code
perfectly. We plan on implementing more optimisations that will eventually
translate the original \hsinl{count} function to

\begin{minted}{haskell}
count'' n' acc' =
  if n' == 0 then acc' else
      count''
          (n' - 1)
          (n' + acc')
count' n acc = count'' (force n) (force acc)
count = delay count'
\end{minted}

Since the body of \hsinl{count''} does not include any delays or any forces, it
can be compiled as-is to CakeML which should then compile it to a loop over the
two integer arguments. It should be noted that, along with the transformations
to the function itself, other optimisations can also occur to the rest of the
program wherever a call to \hsinl{count} appears. These transformations are
challenging to model in a verified compiler, and this is where our work will be
focused.

Lastly, the third part in \autoref{chap:thunks} aims to optimise the
representation of thunks in CakeML. Even with the optimisations discussed for
part two, which try to minimise the generation of unnecessary thunks, many
thunks will eventually make it to the runtime. Currently, CakeML doesn't support
thunks, so PureCake compiles them to CakeML's mutable arrays. Our plan is to add
thunks to the CakeML compilation pipeline so that PureCake can explicitly target
them. This will allow the introduction of runtime optimisations as, for example,
the garbage collector will be able to recognise which data in the heap are
thunks and apply optimisations such as replacing the thunk with its evaluated
value \cite{lazygc}.

\section{Halftime Report: Current status of the Project}

For the first part of the thesis, developing a logical relation for
\ThunkLang{}, the relation is complete and we have proved some theorems about
it, but it's not in a usable state yet. Ideally, we want to use the logical
relation for developing the proof in the second part and then contrast its
implementation with the rest of the proofs that use syntactic relations.
Therefore, we are now working on the third part of the project and have left the
second part for last, hoping this will give us enough time to complete the
relation.

For the third part, we have added thunks to CakeML's source language and adapted
the last intermediate language of PureCake to have explicit thunks and target
CakeML's thunks directly. We are now finishing up on fixing the proofs these
changes broke, and with all that done we will be ready to implement thunk
inlining in the garbage collector.

Work on the second part hasn't started yet, but adding thunks to \StateLang{}
for the third part has demanded changes in \ThunkLang{} too, which has helped
develop a thorough understanding of how \ThunkLang{}'s proofs work and how they
are written. This will make the implementation of a new proof easier to develop
and link up with the rest of the proofs.
